\begin{obeactivity}
\textbf{Aktivitas 2.1: Analisis paket dengan Wireshark}
Gunakan Wireshark untuk menangkap paket data dari lalu lintas yang tidak
terenkripsi. Cari paket yang memuat nama pengguna atau kata sandi dalam bentuk
plainteks. Diskusikan mengapa enkripsi mencegah serangan pasif.
\end{obeactivity}

\begin{obeactivity}
\textbf{Aktivitas 2.2: Mengukur kebocoran waktu pada pembanding kata sandi}
Tulis dua fungsi pemeriksaan kata sandi: satu yang berhenti pada huruf pertama
yang berbeda, satu lagi yang memeriksa seluruh huruf sebelum memutuskan.
\begin{enumerate}
    \item Ukur lama eksekusi untuk kata sandi yang huruf pertamanya sudah salah
          dan yang salah pada huruf terakhir. Berapa selisihnya?
    \item Ulangi pengukuran beberapa ratus kali, lalu hitung rata-ratanya.
    \item Simpulkan mengapa perbedaan waktu sekecil apa pun masih berbahaya.
\end{enumerate}
\end{obeactivity}

\begin{obeactivity}
\textbf{Aktivitas 2.3: Memutar ulang sebuah permintaan}
Buat layanan sederhana yang menerima permintaan berisi perintah dan jumlahnya,
lalu mencatat saldo. Kirim satu permintaan dan simpan isinya.
\begin{enumerate}
    \item Kirim ulang permintaan yang sama persis. Apakah layanan menerimanya?
    \item Tambahkan nomor urut pada permintaan, lalu ubah layanan agar menolak
          nomor yang sudah dipakai. Ulangi pengiriman ulang tadi.
    \item Diskusikan mengapa nomor urut mencegah serangan itu.
\end{enumerate}
\end{obeactivity}
